\section{Discussão e ameaças à validade}

A seleção de benchmarks envolve uma tensão entre atualidade, diversidade e custo de execução. Um conjunto atualizado pode reduzir algumas formas de exposição a tarefas públicas, mas não permite declarar ausência de contaminação sem evidência específica. A cobertura multilíngue, por sua vez, não garante que tarefas de linguagens distintas sejam comparáveis em dificuldade, tamanho ou domínio.

O contraste de localização com contexto mínimo observado por Prathifkumar et al. no SWE-bench Verified reforça essa ameaça \cite{Prathifkumar2025ModelMemory}. Contudo, desempenho superior em um conjunto mais antigo também pode refletir diferenças de distribuição e familiaridade geral com seus repositórios. A interpretação de um eventual contraste entre SWE-bench-Live, Multi-SWE-bench e Terminal-Bench exigirá levar em conta tarefa, linguagem, dificuldade, data e protocolo; uma queda de pontuação em tarefas recentes não será atribuída automaticamente a memorização. Separar desenvolvimento e avaliação por tarefa e, quando viável, por repositório reduz sobreajuste do scaffold às instâncias observadas, mas não demonstra que os modelos nunca tiveram acesso a elas.

Os testes de um benchmark são uma medida operacional de correção e podem não cobrir todos os requisitos de uma issue. Assim, a taxa de resolução deverá ser interpretada como desempenho segundo o protocolo adotado, não como garantia geral de qualidade, segurança ou manutenibilidade do código. A decisão interna do agente de encerrar também deve ser distinguida do veredito externo.

Gorinova et al. argumentam que benchmarks de programação apresentam três desalinhamentos com a engenharia de software agêntica: confundem componentes do sistema em uma pontuação única, podem vincular a correção à forma de uma solução de referência e oferecem pouco sinal sobre onde uma execução falhou \cite{Gorinova2026MisalignedBenchmarks}. Trata-se de um artigo de posição, cujas propostas não estabelecem a frequência desses problemas nos conjuntos aqui selecionados. Para esta pesquisa, a implicação operacional é combinar o veredito oficial com ablações de componentes, classificação de falhas e, em uma amostra delimitada, verificação comportamental independente. Uma solução alternativa que falhe em testes dependentes da implementação de referência e um patch que passe em testes insuficientes representam erros de medição distintos; ambos deverão ser examinados antes de extrapolar a taxa de resolução para qualidade de software em uso.

Modelos e serviços podem mudar ao longo do tempo, e preços, limites de uso e latência podem variar por provedor. A reprodutibilidade exigirá registrar os identificadores e versões dos modelos, os parâmetros, as datas de acesso, os preços e eventuais falhas de serviço. A comparação entre scaffolds também pode ser confundida por diferenças de ferramentas, prompts, orçamento ou harness; esses fatores deverão ser controlados ou documentados.

O estudo controlado em GAIA encontrou efeitos de scaffold dependentes de modelo e dificuldade, inclusive sem a redução monotônica esperada para modelos mais capazes no nível difícil \cite{Starace2026ScaffoldGAIA}. Nesse nível, a vantagem do scaffold com múltiplos papéis apareceu com maior clareza na família Anthropic, mas não nos modelos de outros provedores examinados. Os scaffolds estruturados também fizeram menos chamadas de ferramenta, terminaram mais rapidamente em média e recuperaram-se com maior frequência de erros intermediários; isso mostra que acrescentar etapas de organização não implica necessariamente mais ações ou maior gasto. Ainda assim, sua transferência para engenharia de software é limitada pelo domínio avaliado e pelo fato de o controle ReAct usar uma superfície de ferramentas diferente da dos outros scaffolds. Já Vats e Golev observaram grande variação em tokens por tarefa resolvida entre harnesses de programação, mas a amostra de 50 tarefas e diferenças nos controles de turnos limitam a generalização do fator relatado para custo monetário ou para os modelos futuros \cite{Vats2026ScaffoldEffect}. Portanto, resultados anteriores justificam medir a interação modelo--scaffold e o consumo efetivo, sem fixar sua magnitude no presente estudo.

A redução do conjunto de tarefas introduz outra ameaça. O filtro por dificuldade intermediária estudado por Ndzomga foi eficaz para preservar rankings de agentes em certos conjuntos, mas não assegurou calibração das pontuações absolutas sob mudança de scaffold; além disso, requer resultados históricos suficientes e pode falhar quando poucas tarefas ocupam a faixa selecionada \cite{Ndzomga2026EfficientBenchmarking}. Assim, um ranking obtido na triagem não será usado como evidência de não inferioridade no conjunto completo. A seleção final e sua incerteza deverão ser descritas em termos da população de tarefas que se pretende representar. Se o orçamento permitir apenas um recorte não probabilístico, as conclusões serão restritas a esse recorte.

Por fim, resultados em um conjunto, idioma ou par de modelos não serão generalizados automaticamente a outras linguagens, repositórios ou modelos. As fontes examinadas fundamentam a necessidade de medir qualidade, recursos e validade em conjunto, mas não estabelecem que o roteamento proposto seja superior a agentes de modelo único. Essa conclusão dependerá de uma comparação pareada sob orçamento documentado e de uma margem de qualidade definida previamente. Se os critérios de qualidade, custo e latência apontarem em direções distintas, o resultado deverá ser apresentado como uma troca entre dimensões, e não como superioridade geral.
